\subsection{Calibration and the finite transverse record}
\label{subsec:v47-intro-calibration}

The quantitative inverse also admits a physical acquisition in estimated
contact frames.  There are three distinct approximation errors.  A bridge
with $j$ flights contributes its relative-law error.  Programming its time
using an estimated gap contributes the amplified error
$j|\widehat g-g|$.  Replacing an exact transverse coordinate by one with
error $r$ contributes the probability of crossing a grid edge, including
an outer edge of the recorded square.  These last events are not bounded
by applying a Lipschitz test inequality to a cell indicator.

Theorem~\ref{thm:v47-calibrated-confidence} propagates the three errors
through the same law-to-table modulus, with explicit confidence and
preparation bounds.  Corollary~\ref{cor:v47-physical-histograms} gives a
finite joint prescription for the grid, the final even flight number,
the resolution of a pilot using that same flight number, and the number
of successful categories.  The post-pilot reconstruction retains only
the estimated gaps and transverse categories.  The richer physical
positions used by the pilot and the observable projection remain part
of its acquisition specification; they are not silently inserted into
the compressed inverse.

The structural input throughout is the nonlinear relative profile and
its signed finite-remainder inverse.  Offset normalization, categorical
coupling, concentration, and analytic continuation propagate this input
to different observation levels.  Their combination yields the stated
geometric and statistical consequences without identifying those
observation levels with each other.
